\section{Dasar Pemilihan Atribut Kebijakan ABAC}

Pemilihan atribut kebijakan ABAC dalam penelitian ini didasarkan pada kebutuhan untuk merepresentasikan keputusan akses IoT secara lebih komprehensif, bukan hanya berdasarkan identitas subjek. Secara konseptual, ABAC menentukan keputusan akses berdasarkan atribut subjek, objek, aksi, dan lingkungan, sehingga atribut yang dipilih perlu mencakup identitas pemohon, karakteristik sumber daya, jenis aksi, serta kondisi lingkungan akses \parencite{hu_guide_2014}. Oleh karena itu, atribut $role$ digunakan sebagai representasi \textit{privilige} dasar subjek, tetapi tidak dijadikan satu-satunya dasar keputusan agar model tetap mencerminkan karakteristik ABAC yang \textit{fine-grained} dan \textit{context-aware}. \textcite{diaz_authorization_2025} menegaskan bahwa otorisasi pada IoT menentukan aksi spesifik yang boleh dilakukan setelah autentikasi, sehingga keputusan akses perlu mempertimbangkan atribut yang lebih luas daripada sekadar identitas atau peran.

Kelompok atribut $trust$ dan $anomaly$ dipilih untuk merepresentasikan tingkat kepercayaan serta indikasi perilaku abnormal pada perangkat atau pengguna. \textcite{waheed_blockchain-based_2025} menunjukkan bahwa \textit{dynamic} ABAC pada lingkungan \textit{smart energy} berbasis IoT memanfaatkan mekanisme \textit{trust recalibration} dan analisis perilaku untuk menyesuaikan keputusan akses secara adaptif dan kontekstual. Sementara itu, \textcite{mao_zero-trust_2025} mengintegrasikan \textit{dynamic trust evaluation, anomaly traffic detection,} dan \textit{abnormal log detection} sebagai bagian dari perhitungan nilai risiko dalam model \textit{zero-trust access control} berbasis atribut. Dengan demikian, atribut $trust$ dan $anomaly$ relevan digunakan karena keduanya bersifat dinamis dan mampu merepresentasikan perubahan perilaku aktual entitas selama proses akses berlangsung.

Atribut $network$ dan $pipAvailable$ dipilih untuk merepresentasikan kondisi lingkungan jaringan serta ketersediaan sumber atribut dalam proses evaluasi kebijakan. \textcite{kalaria_adaptive_2024} menunjukkan bahwa \textit{adaptive context-aware access control} pada IoT perlu mempertimbangkan konteks yang berubah, termasuk \textit{network context} seperti kondisi jaringan, trafik jaringan, dan komunikasi antar-\textit{node}, agar keputusan akses dapat disesuaikan secara dinamis. Dalam konteks \textit{edge-IAM}, \textcite{wang_edge-enabled_2025} menunjukkan bahwa LAA diperlukan untuk mendukung akses berlatensi rendah dan tetap berjalan ketika konektivitas \textit{edge-cloud} tidak selalu tersedia dengan memanfaatkan data IAM lokal, \textit{attribtue freshness}, dan sinkronisasi berbasis konteks. Oleh karena itu, $pipAvailable$ dapat digunakan sebagai abstraksi eksperimental untuk menguji kondisi ketika sumber atribut atau layanan informasi kebijakan tidak tersedia, sedangkan $network$ merepresentasikankondisi konektivitas yang memengaruhi kemampuan sistem memperoleh atribut terbaru.

Atribut $firmware$ digunakan untuk merepresentasikan postur keamanan perangkat. \textcite{bakhshi_review_2024} menunjukkan bahwa \textit{firmware} pada perangkat IoT merupakan komponen kritis yang sering terabaikan, karena kerentanan \textit{firmware} dapat menjadi target serangan dan perangkat yang tidak dapat diperbarui akan tetap rentan terhadap kelemahan yang ditemukan setelah perangkat dirilis. Selain itu, aspek \textit{image management, storage integrity,} dan \textit{update delivery} juga berpengaruh terhadap integritas \textit{firmware}, terutama ketika perangkat menjalankan \textit{firmware} usang atau proses pembaruannya tidak aman. \textcite{altaibek_survey_2025} juga menegaskan bahwa \textit{firmware backdoor, command-and-control} tersembunyi, dan manipulasi sensor merupakan \textit{attack vector} yang dapat merusak integritas \textit{deployment} IoT. Dengan demikian, meskipun $firmware$ relatif lebih stabil dibanding atribut perilaku seperti $trust$ atau $anomaly$, status \textit{firmware} tetap relevan digunakan sebagai atribut keamanan perangkat karena \textit{firmware} yang usang, tidak aman, atau tidak tervalidasi dapat meningkatkan risiko kompromi dan seharusnya memengaruhi keputusan akses.

Atribut $objectType$, $action$, dan $sensitivity$ dipilih agar model kebijakan merepresentasikan ABAC secara utuh, bukan sekadar \textit{subject-based access control}. NIST SP 800-162 menjelaskan bahwa otorisasi dalam ABAC ditentukan dengan mengevaluasi atribut yang terkait dengan subjek, objek, operasi yang diminta, dan dalam beberapa kasus kondisi lingkungan terhadap kebijakan atau aturan akses \parencite{hu_guide_2014}. \textcite{ullah_survey_2023} juga menjelaskan bahwa ABAC mendukung kontrol akses yang dinamis, \textit{fine-grained}, dan \textit{context-aware} dengan mempertimbangkan atribut subjek, sumber daya/objek, serta kondisi lingkungan dalam lingkungan IoT. Oleh karena itu, $objectType$ digunakan untuk membedakan jenis sumber daya, $action$ untuk merepresentasikan operasi akses seperti $read$, $write$, atau $control$, sedangkan $sensitivity$ digunakan sebagai atribut objek yang menunjukkan tingkat kritikalitas atau risiko sumber daya yang diakses.

Atribut $critical$ dan $revoked$ dipilih untuk mendukung evaluasi kebijakan yang peka terhadap risiko, terutama pada skenario akses berdampak tinggi, pencabutan hak akses, dan potensi \textit{false allow} akibat \textit{stale attribute}. \textcite{wang_edge-enabled_2025} menunjukkan bahwa keputusan akses pada \textit{edge-IAM} dapat menjadi tidak andal apabila sistem mengasumsikan nilai atribut selalu mutakhir, karena data yang telah diperbarui atau dicabut di \textit{cloud} dapat tetap digunakan secara lokal di \textit{edge} akibat keterlambatan atau kegagalan sinkronisasi. \textcite{azad_verify_2024} juga menekankan bahwa \textit{zero-trust} membutuhkan verifikasi dan evaluasi akses secara kontinu berdasarkan faktor seperti identitas, kondisi perangkat, lokasi, perilaku, dan prinsip \textit{least-privilege}. Dengan demikian, $revoked$ relevan untuk menguji risiko \textit{revocation delay} dan penggunaan atribut yang sudah tidak valid, sedangkan $critical$ digunakan sebagai abstraksi eksperimental untuk membedakan permintaan akses berdampak tinggi yang memerlukan evaluasi lebih ketat.

Berdasarkan karakteristik tersebut, setiap atribut memiliki tingkat volatilitas dan dampak keamanan yang berbeda. Atribut seperti $trust$, $anomaly$, dan $network$ cenderung berubah cepat sehingga membutuhkan \textit{freshness} yang lebih tinggi. Sebaliknya, atribut seperti $role$ dan $firmware$ relatif lebih stabil, tetapi dapat menimbulkan risiko keamanan serius apabila nilai yang usang tetap digunakan dalam keputusan akses. Perbedaan ini menunjukkan bahwa strategi TTL dan \textit{cache} atribut tidak seharusnya diterapkan secara seragam, melainkan perlu disesuaikan dengan volatilitas atribut, sensitivitas sumber daya, serta konsekuensi keamanan dari keputusan akses yang salah.
